\documentclass[10pt,a4paper]{amsart}
\usepackage[margin=19mm]{geometry}
\usepackage{amsmath,amssymb,mathtools}
\usepackage[scr=rsfs,cal=euler]{mathalfa}
\usepackage{xfrac}
\usepackage[unicode,hidelinks,bookmarks=false]{hyperref}
\newcommand{\ie}{i.e.\ }
\newcommand{\cf}{cf.\ }
\newcommand{\resp}{respectively\ }
\newcommand{\Subsection}{\subsection}
\providecommand{\nobreakdash}{\nobreak\hbox{-}}
\setlength{\emergencystretch}{2em}
\multlinegap0pt
\usepackage{xspace,mathrsfs}
\newtheorem{theorem}{Theorem}[section]
\newtheorem{lemma}[theorem]{Lemma}
\newtheorem{definition}[theorem]{Definition}
\newtheorem{example}[theorem]{Example}
\newtheorem{exercise}[theorem]{Exercise}
\newtheorem{proposition}[theorem]{Proposition}
\newtheorem{corollary}[theorem]{Corollary}
\newtheorem{lem}[theorem]{Lemma}
\newtheorem{note}[theorem]{Note}
\newtheorem{remark}[theorem]{Remark}
\numberwithin{equation}{section}
\newcommand{\abs}[1]{\LVERT#1\RVERT}
\begin{document}
\title[(n,Q)-ideals ans phi-(n,Q)-ideals of commutative rings]{(n,Q)-ideals ans phi-(n,Q)-ideals of commutative rings}
\author{[2020] 13A15, 13E99 ęywords $(n,Q)$-ideal, $ $-$(n,Q)$-ideal; Mahdi Anbarloei}
\date{}
\maketitle
\noindent\emph{Source and licence.} (n,Q)-ideals ans phi-(n,Q)-ideals of commutative rings. Version arXiv:2606.14576v1 (CC BY 4.0). This material is distributed under the Creative Commons Attribution 4.0 International licence, http://creativecommons.org/licenses/by/4.0/, as recorded in the arXiv OAI licence metadata for this exact version and shown on the abstract page https://arxiv.org/abs/2606.14576v1. Full rights notice: licenses/arxiv-2606.14576-CC-BY-4.0.txt.\par\smallskip
\noindent\textit{Editorial scope.} Counted unit: the original \texttt{theorem} environment \texttt{6} at source lines 213--219 together with its complete proof at lines 220--267; the delivered document keeps source lines 113--268 so that every referenced label and every citation resolves inside it. Native editable source is the paper's own concatenated TeX; the AMS wrapper is editorial formatting only. No publisher class, upstream script or unrelated artwork is redistributed.
\section{$(n,Q)$-ideals}

\begin{definition}
Let $I$ and $ Q$  be proper ideals of $A$. We say that $I$ is  an  $(n,Q)$-ideal of $A$ if whenever $a_1\cdots a_{n+1} \in I  $ for $a_1,\ldots,a_{n+1} \in A$, then $a_1\cdots  a_n \in I$ or there exists $1 \leq t \leq n$ such that $a_1\cdots \widehat{a_t} \cdots a_{n+1} \in Q.$
\end{definition}
\begin{example}
Let $A=\mathbb{Z}[x,y,z]$, $Q=\langle xy,z  \rangle$ and  $I=\langle xy \rangle$. Then, $I$  is a $(2,Q)$-ideal of $A$, but it is not $Q$-ideal. Indeed, $x  y \in I$ and $x,y \notin Q$ but  $y,x \notin I$.
\end{example}
\begin{proposition} \label{1}
Let $I $   be a proper ideal of $A$. If $I$ is an  $(n,Q)$-ideal of $A$ where $Q$ is a proper ideal of $A$, then $I   \subseteq Q$.
\end{proposition}
\begin{proof}
Suppose that $I$ is an $(n,Q)$-ideal for some proper ideal $Q$ of $A$  such that  $I    \nsubseteq Q$. Then, there exists $a \in I$ such that $a \notin Q$. Since $I$ is a  $(n,Q)$-ideal of $A$ and $1 \cdots 1\cdot  a\in I$ , we get $1  \in I$ which is impossible. Hence, $I \subseteq Q$, as needed.
\end{proof}
Recall from \cite{Khalfi} that a proper ideal $I$ of $A$ is an  $(n,J)$-ideal if whenever $a_1\cdots a_{n+1} \in I  $ for $a_1,\ldots,a_{n+1} \in A$, then $a_1\cdots  a_n \in I$ or there exists $1 \leq t \leq n$ such that $a_1\cdots \widehat{a_t} \cdots a_{n+1} \in J.$ 
\begin{theorem}\label{2}
Let $I$ be a proper ideal of $A$. Then, $I$ is an $(n,J)$-ideal of $A$ if and only if $I$ is an $(n,M)$-ideal for each  $M \in \text{Max}(A)$.
\end{theorem}

\begin{proof}
Let $I$ be a $n$-$J$-ideal of $A$ and $M \in \text{Max}(A)$. Suppose that $a_1\cdots a_{n+1} \in I$ for $a_1,\ldots,a_n \in A$ and $a_1\cdots \widehat{a_t} \cdots a_{n+1} \notin M$ for some $1 \leq t \leq n$. It follows that  $a_1\cdots \widehat{a_t} \cdots a_{n+1} \notin \text{J}(A)$. Since $I$ is an $(n,J)$-ideal of $A$, we get $a_1\cdots a_n \in I$. Thus, $I$ is a $(n,M)$-ideal of $A$. Conversely, assume that  $I$ is a $n$-$M$-ideal for each  $M \in \text{Max}(A)$. So, we obtain  $I \subseteq M$ for every $M \in \text{Max}(A)$ by Proposition \ref{1}. This  means $\text{J}(A)$ contains $I$. Now, suppose that $a_1\cdots a_{n+1} \in I$ for $a_1,\ldots,a_n \in A$ and $a_1\cdots \widehat{a_t} \cdots a_{n+1} \notin \text{J}(A)$ for some $1 \leq t \leq n$. Then, we conclude that  $a_1\cdots \widehat{a_t} \cdots a_{n+1} \notin M_0$ for some $M_0 \in \text{Max}(A)$. By the hypothesis, $I$  is an $(n,M_0)$-ideal of $A$. Therefore, we get $a_1\cdots a_n \in I$. Thus,  $I$ is an $(n,J)$-ideal of $A$.
\end{proof}
\begin{proposition} \label{2.1}
Assume that  $A$ is a commutative ring. 
\begin{enumerate}
\item If $I$ is an  $(n,Q)$-ideal   of $A$ for each ideal $Q$   properly containing $I$, then $I$ is an $n$-absorbing primary ideal.
\item  If $I$ is an $(n,Q)$-ideal of $A$ for each ideal $Q$   properly containing $\sqrt{I}$, then $\sqrt{I}$ is an $n$-absorbing  ideal.
\item The following conditions are equivalent:
\begin{enumerate}
\item[(i)] $I$ is an  $n$-absorbing primary ideal.
\item[(ii)] $I$ is an $(n,\sqrt{I})$-ideal.
\item[(iii)] $I$ is an $(n,Q)$-ideal of $A$ for each ideal $Q$ of $A$ containing $\sqrt{I}$.
\end{enumerate}
\end{enumerate}
\begin{proof}
(1) Suppose that $a_1\cdots a_{n+1} \in I$ for $a_1,\ldots, a_{n+1} \in A$ but $a_1 \cdots a_n \notin I$. Set, $Q=I+\langle a_1 \cdots a_n \rangle $. Then, $I$ is an $(n,Q)$-ideal of $A$ as $I \subsetneq Q$. This implies that $a_1 \cdots \widehat{a_t} \cdots a_{n+	1} \in Q$ for some $1 \leq t \leq n$. Then, there exist $x \in I$ and $y \in A$ such that $a_1 \cdots \widehat{a_t} \cdots a_{n+	1}=x+y(a_1 \cdots a_n)$. Hence, \\

$\hspace{0.5cm}(a_1 \cdots \widehat{a_t} \cdots a_{n+	1})^2=x(a_1 \cdots \widehat{a_t} \cdots a_{n+	1})+y(a_1 \cdots a_n)(a_1 \cdots \widehat{a_t} \cdots a_{n+	1})$

$\hspace{3.4cm}=x(a_1 \cdots \widehat{a_t} \cdots a_{n+	1})+y(a_1 \cdots a_{n+1})(a_1 \cdots \widehat{a_t} \cdots a_n)$.

$\hspace{3.4cm} \in I,$\\

which implies $ a_1 \cdots \widehat{a_t} \cdots a_{n+	1} \in \sqrt{I}$. Therefore, $I$ is an $n$-absorbing primary ideal.

(2) Let $I$ be  an $(n,Q)$-ideal of $A$ for each ideal $Q$ of $A$ properly containing $\sqrt{I}$. Suppose  that $a_1 \cdots a_{n+1} \in \sqrt{I}$ and $a_1 \cdots a_n \notin \sqrt{I}$. This means  that $(a_1 \cdots a_{n+1})^k=a_1^k \cdots a_{n+1}^k\in I$ for some $k \in \mathbb{N}$. Put $Q=\sqrt{I}+\langle a_1 \cdots a_n \rangle$. By the hypothesis, $I$ is an $(n,Q)$-ideal as $\sqrt{I} \subsetneq Q$. Since $I$ is  an $(n,Q)$-ideal of $A$, $a_1^k \cdots a_{n+1}^k\in I$ and $a_1^k \cdots a_n^k\notin I$, we obtain $(a_1 \cdots \widehat{a_t} \cdots a_{n+1})^k=a_1^k \cdots \widehat{a_t^k} \cdots a_{n+1}^k\in Q$ for some $1 \leq t \leq n$. There exists $x \in \sqrt{I}$ and $y \in A$ such that $(a_1 \cdots \widehat{a_t} \cdots a_{n+1})^k=x+y(a_1 \cdots a_n)$. Therefore, we have\\


$(a_1 \cdots \widehat{a_t} \cdots a_{n+1})^{2k}=x(a_1 \cdots \widehat{a_t} \cdots a_{n+1})^k+y(a_1 \cdots a_n)(a_1 \cdots \widehat{a_t} \cdots a_{n+1})^k$

$\hspace{3.1cm}=x(a_1 \cdots \widehat{a_t} \cdots a_{n+1})^k+y(a_1 \cdots a_{n+1})(a_1 \cdots \widehat{a_t} \cdots a_n)^k a_{n+1}^{k-1}$


$\hspace{3.1cm} \in \sqrt{I}.$\\
Hence, $a_1 \cdots \widehat{a_t} \cdots a_{n+1} \in\sqrt{I}$. Consequently, $\sqrt{I}$ is an $n$-absorbing  ideal.

(3) (i) $\Longrightarrow$ (ii) and (ii) $\Longrightarrow$ (iii) are clear.


(iii) $\Longrightarrow$ (i)  Assume that $a_1 \cdots a_{n+1} \in I$ for $a_1,\ldots, a_{n+1} \in A$. Since $I$ is an $(n,\sqrt{I})$-ideal by the hypothesis, we have $a_1\cdots a_n \in I$ or $a_1 \cdots \widehat{a_t} \cdots a_{n+1} \in \sqrt{I}$ for some $1 \leq t \leq n$, as desired. 
\end{proof}
\end{proposition}
\begin{theorem} \label{3}
Let $I \subseteq Q$ be ideals of $A$. If $Q$ is an $(n-1)$-absorbing ideal of $A$, then $I$ is an $(n,Q)$-ideal of $A$. 
\end{theorem}
\begin{proof}
Assume that $a_1\cdots a_{n+1} \in I$ for some $a_1,\ldots,a_{n+1} \in A$ but $a_1\cdots a_n \notin I$. Let $Q$ be  an $(n-1)$-absorbing ideal of $A$. If $a_1\cdots a_n \in Q$, then we obtain $a_1\cdots a_{n-1} \in Q$ or $a_1\cdots \widehat{a_t} \cdots a_{n} \in Q$ for some $1 \leq t \leq n-1$. This implies that $a_1\cdots a_{n-1} \widehat{a_n} a_{n+1} \in Q$ or $a_1\cdots \widehat{a_t} \cdots a_{n}a_{n+1} \in Q$ for some $1 \leq t \leq n-1$.
If $a_1\cdots a_n \notin Q$, then we have $(a_1a_2)a_3\cdots  \widehat{a_t}\cdots a_{n+1} \in Q$ for some $1 \leq t \leq n$ as $Q$ is an $(n-1)$-absorbing ideal of $A$ and $(a_1a_2)a_3\cdots   a_{n+1} \in Q$. Thus, $I$ is an $(n,Q)$-ideal of $A$. 
\end{proof}



\begin{proposition} \label{4}
Let $Q$ be an   ideal of $A$. If $I_j$ is an $(n_j,Q)$-ideal of $A$ for each $1 \leq j \leq k$, then $\cap_{j=1}^k I_j$ is an $(n,Q)$-ideal where $n=\Sigma_{j=1}^k n_j$. 
\end{proposition}
\begin{proof}
Let $a_1 \cdots a_{n+1} \in \cap_{j=1}^k I_j$ for $a_1,\ldots,a_{n+1} \in A$ but $a_1\cdots \widehat{a_t} \cdots a_{n+1} \notin Q$ for all $1 \leq t \leq n$.  Since $I_j$ is an $(n_j,Q)$-ideal of $A$ and $a_1 \cdots a_{n+1} \in I_j$ for all $1 \leq j \leq k$, there exist elements  $1 \leq \alpha_1^j, \alpha_2^j,\ldots,\alpha_{n_j}^j \leq n$ with $a_{\alpha_1^j}, a_{\alpha_2^j},\ldots,a_{\alpha_{n_j}^j}  \in I_j$. Let $\alpha_l^r=\alpha_m^s$ for pairs $r,l$ and $s,m$. Then,  


\[a_{\alpha_1^1}  a_{\alpha_2^1}  \cdots  a_{\alpha_{n_1}^1}   \cdots a_{\alpha_1^r}  a_{\alpha_2^r}   \cdots   a_{\alpha_l^r}   \cdots   a_{\alpha_ {n_r}^r}   \cdots    a_{\alpha_1^s}   a_{\alpha_2^s}\]
\[ \cdots       \widehat{a_{\alpha_m^s}} \cdots  a_{\alpha_{n_s}^s} \cdots    a_{\alpha_1^k}   a_{\alpha_2^k}   \cdots   a_{\alpha_{n_k}^k} \in Q \]
which implies $a_1 \cdots \widehat{a_{\alpha_m^s}} \cdots a_na_{n+1} \in Q$ which is a contradiction. Therefore,  ${\alpha_i^j}^,$s are different. From $a_{\alpha_1^j}, a_{\alpha_2^j},\ldots,a_{\alpha_{n_j}^j}  \in I_j$ for all $1 \leq j \leq n$ it follows that $a_1 \cdots a_n=a_{\alpha_1^1}  a_{\alpha_2^1}  \cdots  a_{\alpha_{n_1}^1} \cdots a_{\alpha_1^k}   a_{\alpha_2^k}   \cdots   a_{\alpha_{n_k}^k} \in \cap_{j=1}^k I_j$. Consequently, $\cap_{j=1}^k I_j$ is an $(n,Q)$-ideal of $A$.
\end{proof}

Let     $I$ be a proper ideal of $A$.  Then, we  use $Z_I(A)$ to denote  the set $\{x \in A \mid xs \in I \  \text{for some} \ s \in A \backslash I\}$.
 \begin{theorem} \label{5}
Let $I$ and $Q$ be ideals of $A$ and $S$ be a multiplicative subset of $A$. 
\begin{enumerate}
\item If $I$ is an $(n,Q)$-ideal of $A$ such that  $Q \cap S=\varnothing$, then $S^{-1}I$ is an $(n,S^{-1}Q)$-ideal of $S^{-1}A$.
\item If $S^{-1}I$ is  an $(n,S^{-1}Q)$-ideal of $S^{-1}A$ and $S \cap Z_j=\varnothing$ for $j \in\{I,Q\}$, then $I$ is an $(n,Q)$-ideal of $A$.
\end{enumerate} 
\end{theorem}
\begin{proof}
(1) Let $I$ be an $(n,Q)$-ideal of $A$. Assume that $  \frac{a
_1}{s_1}\cdots\frac{a_{n+1}}{s_{n+1}} \in S^{-1}I$ for  $ \frac{a
_1}{s_1},\ldots,\frac{a_{n+1}}{s_{n+1}} \in S^{-1}A$ such that $  \frac{a_1}{s_1}\cdots\frac{a_n}{s_n} \notin S^{-1}I$. Hence, we get $ sa_1\cdots a_{n+1}=a_1\cdots a_n(sa_{n+1}) \in I$ for some $s \in S$. Since $I$ is an $(n,Q)$-ideal  of $A$ and $a_1\cdots a_n \notin I$, we get $a_1\cdots \widehat{a_t} \cdots a_n (sa_{n+1})\in Q$ for some $1 \leq t \leq n$ which means  $ \frac{a
_1}{s_1}\cdots \widehat{\frac{a_t}{s_t}} \cdots \frac{a_{n+1}}{s_{n+1}} \in S^{-1}Q$. Thus,  $S^{-1}I$ is an $(n,S^{-1}Q)$-ideal of $S^{-1}A$.

(2) Let  $S^{-1}I$ be an $(n,S^{-1}Q)$-ideal of $S^{-1}A$. Suppose that $  a_1 \cdots a_{n+1} \in I$ for $a_1, \ldots, a_{n+1} \in A$ and $  a_1 \cdots a_n \notin I$ . This means that $\frac{a_1}{1} \cdots \frac{a_{n+1}}{1} \in S^{-1}I$. If $\frac{a_1}{1} \cdots \frac{a_n}{1} \in S^{-1}I$, then we have $sa_1\cdots a_n \in I$ for some $s \in S$. Since  $S \cap Z_I(A)=\varnothing$, we have  $a_1\cdots a_n \in I$, a contradiction. Since $S^{-1}I$ is an $(n,S^{-1}Q)$-ideal of $S^{-1}A$ and $\frac{a_1}{1} \cdots \frac{a_n}{1} \notin S^{-1}I$, we have $\frac{a_1}{1} \cdots \widehat{\frac{a_t}{1}} \cdots  \frac{a_{n+1}}{1} \in S^{-1}Q$ for some $1 \leq t \leq n$. It follows that $sa_1\cdots \widehat{a_t} \cdots a_{n+1} \in Q$. Since $S \cap Z_Q=\varnothing$, we get $a_1\cdots \widehat{a_t} \cdots a_{n+1} \in Q$. Consequently, $I$ is an $(n,Q)$-ideal of $A$.
\end{proof}

\begin{theorem} \label{6}
Let $A= A_1 \times \cdots \times A_r$ be a decomposable ring and let \[ I = I_1 \times \cdots \times I_{\beta_1-1} \times A_{\beta_1} \times I_{\beta_1+1} \times \cdots \times I_{\beta_k-1} \times A_{\beta_k} \times I_{\beta_k+1} \times \cdots \times I_r \] and \[ Q = Q_1 \times \cdots \times Q_{\beta_1-1} \times A_{\beta_1} \times Q_{\beta_1+1} \times \cdots \times Q_{\beta_k-1} \times A_{\beta_k} \times Q_{\beta_k+1} \times \cdots \times Q_r \]be   ideals of $A$ with    $\{\beta_1, \dots, \beta_k\} \subset \{1, \dots, r\}$.  The following conditions are equivalent:
 \begin{enumerate} 
 \item[(1)] $I$ is an $(n,Q)$-ideal of $A$; 
 \item[(2)] $I'  = I_1 \times \cdots \times I_{\beta_1-1} \times I_{\beta_1+1} \times \cdots \times I_{\beta_k-1} \times I_{\beta_k+1} \times \cdots \times I_r$ is an $(n,Q')$-ideal of $A' =A_1 \times \cdots \times A_{\beta_1-1} \times A_{\beta_1+1} \times \cdots \times A_{\beta_k-1} \times A_{\beta_k+1} \times \cdots \times A_r$ where $Q'=Q_1 \times \cdots \times Q_{\beta_1-1} \times Q_{\beta_1+1} \times \cdots \times Q_{\beta_k-1} \times Q_{\beta_k+1} \times \cdots \times I_r$.
  \end{enumerate}
\end{theorem}

\begin{proof}
(1) $\Longrightarrow$ (2) Let $I$ be an $(n,Q)$-ideal of $A$ and \\

$\hspace{1cm} (a^{(1)}_1,\ldots,a^{(1)}_{\beta_1-1}, a^{(1)}_{\beta_1+1},\cdots, a^{(1)}_{\beta_k-1}, a^{(1)}_{\beta_k+1}, \dots,a_r^{(1)})\cdots$\\

$\hspace{2cm}(a^{(n+1)}_1,\ldots,a^{(n+1)}_{\beta_1-1}, a^{(n+1)}_{\beta_1+1},\ldots, a^{n+1}_{\beta_k-1}, a^{(n+1)}_{\beta_k+1},\ldots,a_r^{(n+1)}) \in I'$

such that  $a_i^{(k),}s$  are in $A_i$. This implies that\\

$\hspace{1cm} (a^{(1)}_1,\ldots,a^{(1)}_{\beta_1-1}, 1, a^{(1)}_{\beta_1+1},\cdots, a^{(1)}_{\beta_k-1}, 1,a^{(1)}_{\beta_k+1}, \dots,a_r^{(1)})\cdots$\\

$\hspace{2cm}(a^{(n+1)}_1,\ldots,a^{(n+1)}_{\beta_1-1},1, a^{(n+1)}_{\beta_1+1},\ldots, a^{(n+1)}_{\beta_k-1},1, a^{(n+1)}_{\beta_k+1},\ldots,a_r^{(n+1)}) \in I.$\\

Since $I$ is an $(n,Q)$-ideal of $A$, we conclude that\\

$\hspace{1cm} (a^{(1)}_1,\ldots,a^{(1)}_{\beta_1-1}, 1, a^{(1)}_{\beta_1+1},\cdots, a^{(1)}_{\beta_k-1}, 1,a^{(1)}_{\beta_k+1}, \dots,a_r^{(1)})\cdots$\\

$\hspace{2cm}(a^{(n)}_1,\ldots,a^{(n)}_{\beta_1-1},1, a^{(n)}_{\beta_1+1},\ldots, a^{(n)}_{\beta_k-1},1, a^{(n)}_{\beta_k+1},\ldots,a_r^{(n)}) \in I$\\

or\\

$\hspace{0.75cm} (a^{(1)}_1,\ldots,a^{(1)}_{\beta_1-1}, 1, a^{(1)}_{\beta_1+1},\cdots, a^{(1)}_{\beta_k-1}, 1,a^{(1)}_{\beta_k+1}, \dots,a_r^{(1)})\cdots$\\

$\hspace{1.75cm}(a^{(t-1)}_1,\ldots,a^{(t-1)}_{\beta_1-1},1, a^{(t-1)}_{\beta_1+1},\ldots, a^{(t-1)}_{\beta_k-1},1, a^{(t-1)}_{\beta_k+1},\ldots,a_r^{(t-1)}) $\\

$\hspace{2cm}(a^{(t+1)}_1,\ldots,a^{(t+1)}_{\beta_1-1},1, a^{(t+1)}_{\beta_1+1},\ldots, a^{(t+1)}_{\beta_k-1},1, a^{(t+1)}_{\beta_k+1},\ldots,a_r^{(t+1)}) \cdots$\\

$\hspace{2.25cm}(a^{(n+1)}_1,\ldots,a^{(n+1)}_{\beta_1-1},1, a^{(n+1)}_{\beta_1+1},\ldots, a^{(n+1)}_{\beta_k-1},1, a^{(n+1)}_{\beta_k+1},\ldots,a_r^{(n+1)}) \in Q$\\

for some $1 \leq t \leq n$. Then, we obtain \\

$\hspace{1cm} (a^{(1)}_1,\ldots,a^{(1)}_{\beta_1-1},  a^{(1)}_{\beta_1+1},\cdots, a^{(1)}_{\beta_k-1}, a^{(1)}_{\beta_k+1}, \dots,a_r^{(1)})\cdots$\\

$\hspace{2cm}(a^{(n)}_1,\ldots,a^{(n)}_{\beta_1-1},  a^{(n)}_{\beta_1+1},\ldots, a^{(n)}_{\beta_k-1},  a^{(n)}_{\beta_k+1},\ldots,a_r^{(n)}) \in I'$\\

or\\

$\hspace{0.75cm} (a^{(1)}_1,\ldots,a^{(1)}_{\beta_1-1},   a^{(1)}_{\beta_1+1},\cdots, a^{(1)}_{\beta_k-1},  a^{(1)}_{\beta_k+1}, \dots,a_r^{(1)})\cdots$\\

$\hspace{1.75cm}(a^{(t-1)}_1,\ldots,a^{(t-1)}_{\beta_1-1},  a^{(t-1)}_{\beta_1+1},\ldots, a^{(t-1)}_{\beta_k-1},  a^{(t-1)}_{\beta_k+1},\ldots,a_r^{(t-1)}) $\\

$\hspace{2cm}(a^{(t+1)}_1,\ldots,a^{(t+1)}_{\beta_1-1},  a^{(t+1)}_{\beta_1+1},\ldots, a^{(t+1)}_{\beta_k-1},  a^{(t+1)}_{\beta_k+1},\ldots,a_r^{(t+1)}) \cdots$\\

$\hspace{2.25cm}(a^{(n+1)}_1,\ldots,a^{(n+1)}_{\beta_1-1},  a^{(n+1)}_{\beta_1+1},\ldots, a^{(n+1)}_{\beta_k-1},  a^{(n+1)}_{\beta_k+1},\ldots,a_r^{(n+1)}) \in Q'$\\
Thus, $I'$ is an $(n,Q')$-ideal of $A'$.

(2) $\Longleftarrow$ (1)   Suppose that $I'  = I_1 \times \cdots \times I_{\beta_1-1} \times I_{\beta_1+1} \times \cdots \times I_{\beta_k-1} \times I_{\beta_k+1} \times \cdots \times I_r$ is an $(n,Q')$-ideal of $A' =A_1 \times \cdots \times A_{\beta_1-1} \times A_{\beta_1+1} \times \cdots \times A_{\beta_k-1} \times A_{\beta_k+1} \times \cdots \times A_r$ where $Q'=Q_1 \times \cdots \times Q_{\beta_1-1} \times Q_{\beta_1+1} \times \cdots \times Q_{\beta_k-1} \times Q_{\beta_k+1} \times \cdots \times I_r$. Using an argument similar to the previous part, we  can   see  that $I$ is an $(n,Q)$-ideal of $A$.
\end{proof}


\begin{thebibliography}{99}
\bibitem[Khalfi]{Khalfi} A.E. Khalfi {\it On $\phi$-$(n,J)$-ideals and $(n,J)$-ideals of commutative rings}, Moroccan Journal of Algebra and Geometry with Applications, {\bf x} (x) (2022) 1-11.
\end{thebibliography}
\end{document}
